% uniud-example-dipartimento.tex
% -----------------------------------------------------------------------------
% ESEMPIO 3 / DOCUMENTAZIONE
% Lettera di dipartimento (qui in modalita' digitale).
%
% L'esempio e' una lettera auto-documentante e descrive in particolare il
% comando \UniudDipartimento{...}, che imposta marchio, acronimo e denominazione
% estesa del dipartimento.
% -----------------------------------------------------------------------------
\documentclass{uniudletter}

\usepackage{polyglossia}
\setmainlanguage{italian}
\usepackage{metalogo}
\usepackage{xcolor}
\usepackage{listings}

\definecolor{UniudCodeComment}{cmyk}{0,0.81,1,0.60}
\definecolor{UniudCodeKeyword}{cmyk}{0.64,0,0.95,0.40}
\definecolor{UniudCodeIdentifier}{cmyk}{0.62,0.57,0.23,0}
\lstset{
  basicstyle=\ttfamily\small,
  keywordstyle=\color{UniudCodeKeyword},
  identifierstyle=\color{UniudCodeIdentifier}\bfseries,
  stringstyle=\color{blue},
  commentstyle=\color{UniudCodeComment},
  morecomment=[l][\color{UniudCodeComment}]{\#},
  breaklines=true,
  columns=fullflexible,
  language=TeX
}

% Il comando imposta automaticamente acronimo e denominazione estesa.
\UniudDipartimento{DPIA}
\UniudDigitale

\UniudSetup{
  luogo = {Udine},
  indirizzo = {via delle Scienze 206},
  cap = 33100,
  citta = Udine,
  provincia = UD,
  paese = Italia,
  telefono = {+39 0432 558400},
  fax = {+39 0432 558499},
  email = {dpia@uniud.it},
  sito = {uniud.it},
  responsabile = {Nome Cognome -- nome.cognome@uniud.it},
  compilatore = {Nome Cognome -- nome.cognome@uniud.it},
  footer-email = false,
  footer-cf-piva = true,
  footer-abi = true,
  footer-cab = true,
  firmatario = {Nome Cognome},
  ruolo = {Direttore del Dipartimento}
}

\UniudDestinatario{
  nome = {Spett.li Utenti \LaTeX},
  struttura = {Universita' degli Studi di Udine},
  indirizzo = {Loro sedi},
  citta = {Udine}
}

\UniudOggetto{Come scrivere una lettera di dipartimento con uniudletter}

\begin{document}
\begin{letter}{}
\opening{Gentili signore e gentili signori,}

per una lettera di dipartimento non e' necessario impostare manualmente
acronimo e denominazione estesa. La classe mette a disposizione un comando
specifico:

\begin{lstlisting}
\documentclass{uniudletter}
\usepackage{polyglossia}
\setmainlanguage{italian}

\UniudDipartimento{DPIA}
\UniudDigitale
\end{lstlisting}

\lstinline|\UniudDipartimento{DPIA}| seleziona l'intestazione di dipartimento,
imposta l'acronimo \texttt{DPIA}, la denominazione estesa e costruisce il
marchio con il sigillo UNIUD e l'acronimo secondo le proporzioni previste dal
manuale di identita' visiva.

Sono riconosciuti direttamente gli acronimi:

\begin{center}
\texttt{DIUM, DILL, DMIF, DPIA, DIES, DISG, DI4A, DMED}.
\end{center}

I dati variabili della sede e della persona che scrive si impostano con
\lstinline|\UniudSetup|:

\begin{lstlisting}
\UniudSetup{
  luogo = {Udine},
  indirizzo = {via delle Scienze 206},
  cap = 33100,
  citta = Udine,
  provincia = UD,
  paese = Italia,
  telefono = {+39 0432 558400},
  email = {dpia@uniud.it},
  sito = {uniud.it},
  firmatario = {Nome Cognome},
  ruolo = {Direttore del Dipartimento}
}
\end{lstlisting}

Le chiavi \texttt{telefono}, \texttt{email} e \texttt{sito} sono proprio i
campi da usare per i riferimenti della lettera. Possono contenere i contatti
generali del dipartimento oppure quelli del singolo mittente, a seconda del
caso. Il numero di fax, se necessario, si imposta con \texttt{fax}.

Per i riferimenti del procedimento sono disponibili anche:

\begin{lstlisting}
\UniudSetup{
  responsabile = {Nome Cognome -- nome.cognome@uniud.it},
  compilatore  = {Nome Cognome -- nome.cognome@uniud.it}
}
\end{lstlisting}

Destinatario e oggetto si indicano esattamente come nelle lettere
istituzionali:

\begin{lstlisting}
\UniudDestinatario{
  nome = {Prof.ssa Nome Cognome},
  struttura = {Universita' di esempio},
  indirizzo = {via Esempio 1},
  cap = 00100,
  citta = Roma,
  provincia = RM,
  paese = Italia
}

\UniudOggetto{Oggetto della comunicazione}
\end{lstlisting}

Infine si scrive normalmente il corpo della lettera:

\begin{lstlisting}
\begin{document}
\begin{letter}{}
\opening{Gentile Professoressa,}

Testo della lettera.

\UniudFirma
\end{letter}
\end{document}
\end{lstlisting}

Nelle lettere di dipartimento il protocollo e i riferimenti sono facoltativi.
Se non viene usato \lstinline|\UniudProtocollo|, il relativo blocco non viene
stampato e non viene riservato spazio. Se invece serve, lo si puo' aggiungere
indipendentemente dalla modalita' scelta.

Questo esempio usa la modalita' digitale. Per ottenere la corrispondente
lettera di dipartimento analogica basta sostituire
\lstinline|\UniudDigitale| con \lstinline|\UniudAnalogico|. Se la lettera
analogica deve riportare il protocollo, si aggiungono i dati:

\begin{lstlisting}
\UniudProtocollo{
  numero = {4},
  titolo = {14},
  classe = {2},
  fascicolo = {1}
}
\end{lstlisting}

Il pie' di pagina della prima pagina e' modulare. Il responsabile e il
compilatore sono facoltativi: se le relative chiavi non vengono valorizzate,
le righe non compaiono. L'indirizzo del footer puo' essere ricavato
automaticamente dai dati gia' usati nell'intestazione; non occorre quindi
scriverlo una seconda volta.

Per default vengono mostrati indirizzo, contatti disponibili e la coppia
CF/P.IVA dell'Ateneo. ABI, CAB, CIN e conto corrente sono disponibili ma
spenti per default. Ogni componente puo' essere attivata o disattivata:

\begin{lstlisting}
\UniudSetup{
  responsabile = {Nome Cognome - nome.cognome@uniud.it},
  compilatore  = {Nome Cognome - nome.cognome@uniud.it},

  footer-indirizzo = true,
  footer-telefono = true,
  footer-fax = false,
  footer-email = false,
  footer-sito = true,
  footer-cf-piva = true,

  footer-abi = true,
  footer-cab = true,
  footer-cin = false,
  footer-conto-corrente = false
}
\end{lstlisting}

I valori istituzionali predefiniti per CF, P.IVA, ABI, CAB, CIN e conto
corrente possono essere sostituiti con le chiavi \texttt{codice-fiscale},
\texttt{partita-iva}, \texttt{abi}, \texttt{cab}, \texttt{cin} e
\texttt{conto-corrente}.

Se si vogliono mantenere i marchi corporate ma nessun dato testuale centrale:

\begin{lstlisting}
\UniudSetup{footer-riferimenti = false}
\end{lstlisting}

Se invece si vuole sopprimere completamente il pie' di pagina della prima
pagina:

\begin{lstlisting}
\UniudSetup{footer-prima-pagina = false}
\end{lstlisting}

La numerazione delle pagine successive non viene disattivata. Marchio e
certificazioni possono inoltre essere controllati separatamente con
\texttt{footer-marchio} e \texttt{footer-certificazioni}. Una nota libera puo'
essere aggiunta con \texttt{pie-prima-pagina}.

Dalla seconda pagina viene usato automaticamente il marchio contratto
\texttt{UNI/UD} largo 21,3 mm, posizionato a 12 mm dai bordi superiore e
sinistro; il numero di pagina e' nel formato \texttt{2 di X}.

Non serve creare un file \texttt{.lco} personalizzato per ciascun dipartimento:
la scelta del dipartimento e i dati variabili della lettera sono separati e
vengono gestiti dall'interfaccia della classe.

\UniudFirma
\end{letter}
\end{document}
